\begin{remark}[2.2.9]
\label{IV.2.2.9-fr}
Los resultados de (2.2.7) y (2.2.8) siguen siendo válidos cuando los
$\mathcal O_Y$-módulos $\mathscr G'$, $\mathscr G$, $\mathscr G''$ que figuran
en ellos son arbitrarios (no necesariamente cuasicoherentes). En efecto,
para todo $y\in Y$, existe $x\in f^{-1}(y)$ tal que $\mathscr F_x$ sea un
$\mathcal O_y$-módulo
\oldpage[IV-2]{12}
fielmente plano y, por consiguiente, el funtor
$\mathscr G_y\leadsto\mathscr G_y\otimes_{\mathcal O_y}\mathscr F_x$ es
fiel; como, por otra parte, para \emph{todo} $x\in f^{-1}(y)$ el funtor
$\mathscr G_y\leadsto\mathscr G_y\otimes_{\mathcal O_y}\mathscr F_x$ es
exacto, se deduce inmediatamente la conclusión.
\end{remark}

\begin{proposition}[2.2.10]
\label{IV.2.2.10-fr}
Sean $f:X\to Y$, $g:Y'\to Y$, $h:Y'\to Z$ tres morfismos de preesquemas,
$X'=X\times_Y Y'$, $\mathscr F$ un $\mathcal O_X$-módulo cuasicoherente,
fielmente plano respecto de $Y$, y $\mathscr G'$ un
$\mathcal O_{Y'}$-módulo cuasicoherente. Para que $\mathscr G'$ sea plano
\textup{(resp.} fielmente plano) respecto de $Z$, es necesario y suficiente que
$\mathscr F\otimes_Y\mathscr G'$ sea un $\mathcal O_{X'}$-módulo plano
\textup{(resp.} fielmente plano) respecto de $Z$.
\end{proposition}

Ya se sabe que, si $\mathscr G'$ es $Z$-plano, lo mismo ocurre con
$\mathscr F\otimes_Y\mathscr G'$ (2.1.5). Consideremos un cambio de base
arbitrario $Z''\to Z$ y pongamos $X''=X'\times_Z Z''$; si $\mathscr G'$ es
fielmente plano respecto de $Z$, el funtor
\begin{equation}
\tag{2.2.10.1}
\label{IV.2.2.10.1-fr}
\mathscr H''\leadsto\mathscr H''\otimes_Z\mathscr G'
\leadsto(\mathscr H''\otimes_Z\mathscr G')\otimes_Y\mathscr F
=\mathscr H''\otimes_Z(\mathscr G'\otimes_Y\mathscr F)
\end{equation}
de la categoría de los $\mathcal O_{Z''}$-módulos cuasicoherentes a la de
los $\mathcal O_{X''}$-módulos es el compuesto de dos funtores exactos y
fieles, por lo que es exacto y fiel. Recíprocamente, si este funtor compuesto
es exacto (resp. exacto y fiel), lo mismo ocurre con
$\mathscr H''\leadsto\mathscr H''\otimes_Z\mathscr G'$, puesto que el funtor
$\mathscr M'\leadsto\mathscr M'\otimes_Y\mathscr F$ (de la categoría de los
$\mathcal O_{Y'}$-módulos cuasicoherentes a la de los
$\mathcal O_X$-módulos cuasicoherentes) es exacto y fiel por hipótesis.

\begin{corollary}[2.2.11]
\label{IV.2.2.11-fr}
\begin{enumerate}
\item[\textup{(i)}] Sean $f:X\to Y$, $g:Y'\to Y$ dos morfismos,
$X'=X\times_Y Y'$, $\mathscr F$ un $\mathcal O_X$-módulo cuasicoherente. Si
$\mathscr F$ es fielmente plano respecto de $Y$,
$\mathscr F'=\mathscr F\otimes_Y\mathcal O_{Y'}$ es fielmente plano
respecto de $Y'$.
\item[\textup{(ii)}] Sean $f:X\to Y$, $g:Y\to Z$ dos morfismos, $\mathscr F$ un
$\mathcal O_X$-módulo cuasicoherente, fielmente plano respecto de $Y$.
Para que $g$ sea un morfismo fielmente plano, es necesario y suficiente que
$\mathscr F$ sea fielmente plano respecto de $Z$.
\item[\textup{(iii)}] Sean $f:X\to Y$, $g:Y\to Z$ dos morfismos, $\mathscr G$ un
$\mathcal O_Y$-módulo cuasicoherente. Se supone que el morfismo $f$ es
fielmente plano. Para que $\mathscr G$ sea plano \textup{(resp.} fielmente plano)
respecto de $Z$, es necesario y suficiente que $f^*(\mathscr G)$ sea plano
\textup{(resp.} fielmente plano) respecto de $Z$.
\item[\textup{(iv)}] Sean $f:X\to Y$, $g:Y\to Z$ dos morfismos, $x$ un punto de
$X$. Se supone que $f$ es plano en el punto $x$. Para que $g\circ f$ sea
plano en el punto $x$, es necesario y suficiente que $g$ sea plano en el
punto $f(x)$.
\end{enumerate}
\end{corollary}

Para demostrar (i), se aplica (2.2.10) reemplazando $Z$ por $Y'$ y
$\mathscr G'$ por $\mathcal O_{Y'}$. Para demostrar (ii), se aplica
(2.2.10) tomando como morfismo $Y'\to Y$ la identidad y reemplazando
$\mathscr G'$ por $\mathcal O_Y$. Para demostrar (iii), se aplica (2.2.10)
tomando de nuevo la identidad para $Y'\to Y$ y reemplazando $\mathscr F$ por
$\mathcal O_X$ y $\mathscr G'$ por $\mathscr G$. Por último, (iv) se deduce
de (ii) reemplazando $X$ por $\operatorname{Spec}(\mathcal O_x)$,
$\mathscr F$ por $\mathcal O_X$, $Y$ por
$\operatorname{Spec}(\mathcal O_{f(x)})$ y $Z$ por
$\operatorname{Spec}(\mathcal O_{g(f(x))})$, y teniendo en cuenta
$(\mathbf{0}_{\mathrm I},6.6.2)$.

\begin{corollary}[2.2.12]
\label{IV.2.2.12-fr}
Sean $Y$ un esquema afín, $f:X\to Y$ un morfismo cuasicompacto,
$\mathscr F$ un $\mathcal O_X$-módulo cuasicoherente. Si $\mathscr F$ es
fielmente plano respecto de $f$, existen un esquema afín $X'$ y un
isomorfismo local sobreyectivo $g:X'\to X$ tales que $g^*(\mathscr F)$ sea
fielmente plano respecto de $f\circ g$.
\end{corollary}

En efecto, $X$ es cuasicompacto, luego es reunión finita de abiertos afines
$X_i$; basta tomar como $X'$ el esquema afín \emph{suma} de los preesquemas
inducidos sobre los abiertos $X_i$ y
\oldpage[IV-2]{13}
como $g:X'\to X$ el morfismo canónico; es claro que $g$ es fielmente plano,
y la hipótesis implica que $g^*(\mathscr F)$ es fielmente plano respecto de
$f\circ g$, en virtud de (2.2.11, (iii)).

\begin{corollary}[2.2.13]
\label{IV.2.2.13-fr}
\begin{enumerate}
\item[\textup{(i)}] Sean $f:X\to Y$, $g:Y'\to Y$ dos morfismos,
$X'=X\times_Y Y'$, $f'=f_{(Y')}:X'\to Y'$. Si $f$ es un morfismo
fielmente plano, lo mismo ocurre con $f'$.
\item[\textup{(ii)}] Si $f:X\to X'$, $g:Y\to Y'$ son dos $S$-morfismos fielmente
planos,
\[
f\times_S g:X\times_S Y\longrightarrow X'\times_S Y'
\]
es fielmente plano.
\item[\textup{(iii)}] Sean $f:X\to Y$, $g:Y\to Z$ dos morfismos tales que $f$ sea
fielmente plano. Para que $g$ sea un morfismo plano \textup{(resp.} fielmente plano),
es necesario y suficiente que $g\circ f$ sea un morfismo plano \textup{(resp.}
fielmente plano).
\end{enumerate}
\end{corollary}

\begin{proposition}[2.2.14]
\label{IV.2.2.14-fr}
Sea $f:X\to Y$ un morfismo fielmente plano y cuasicompacto. Si $X$ es
localmente noetheriano, lo mismo ocurre con $Y$.
\end{proposition}

Como la cuestión es local en $Y$, se puede suponer que
$Y=\operatorname{Spec}(A)$ es afín; como $f$ es cuasicompacto, resulta de
(2.2.12) que también es posible limitarse al caso en que
$X=\operatorname{Spec}(B)$ es afín. Entonces $B$ es un anillo noetheriano
por hipótesis y un $A$-módulo fielmente plano (2.2.3); por tanto, $A$ es
noetheriano $(\mathbf{0}_{\mathrm I},6.5.2)$.

\begin{proposition}[2.2.15]
\label{IV.2.2.15-fr}
Sea $f:X\to Y$ un morfismo fielmente plano. Si $\mathfrak S(X)$ y
$\mathfrak S(Y)$ designan respectivamente el conjunto de los subpreesquemas
de $X$ y de $Y$, la aplicación $Z\leadsto f^{-1}(Z)$ de $\mathfrak S(Y)$ en
$\mathfrak S(X)$ es inyectiva.
\end{proposition}

Como $f$ es sobreyectivo, para el \emph{conjunto} subyacente a un subpreesquema $Z$
de $Y$ se tiene $f(f^{-1}(Z))=Z$. Por otra parte, si $U$ es un abierto de $Y$
que contiene a $Z$ y en el que $Z$ es cerrado, $f^{-1}(U)$ es abierto en $X$,
$f^{-1}(Z)$ es cerrado en $f^{-1}(U)$, y la restricción $f^{-1}(U)\to U$ de
$f$ es un morfismo fielmente plano. Es posible, pues, limitarse a considerar
únicamente subpreesquemas \emph{cerrados} de $Y$. Ahora bien, si $Z$ es un
subpreesquema cerrado de $Y$ correspondiente a un ideal cuasicoherente
$\mathscr J$ de $\mathcal O_Y$, se sabe que $f^{-1}(Z)$ corresponde al ideal
cuasicoherente $f^*(\mathscr J)\mathcal O_X$ de $\mathcal O_X$ (\textbf I, 4.4.5),
y, como $f$ es plano, $f^*(\mathscr J)\mathcal O_X$ se identifica con
$f^*(\mathscr J)$. Pero la aplicación $\mathscr J\leadsto f^*(\mathscr J)$ del
conjunto de los ideales cuasicoherentes de $\mathcal O_Y$ en el conjunto de
los ideales cuasicoherentes de $\mathcal O_X$ es inyectiva (2.2.7), de donde
se sigue la conclusión.

\begin{corollary}[2.2.16]
\label{IV.2.2.16-fr}
Sean $X$, $Y$ dos $S$-preesquemas; si $S'\to S$ es un morfismo fielmente
plano, la aplicación $f\leadsto f_{(S')}$ de $\operatorname{Hom}_S(X,Y)$ en
$\operatorname{Hom}_{S'}(X_{(S')},Y_{(S')})$ es inyectiva.
\end{corollary}

Se tiene $X_{(S')}\times_{S'}Y_{(S')}=(X\times_S Y)_{(S')}$ (\textbf I, 3.3.10), por
lo que el morfismo de proyección
$X_{(S')}\times_{S'}Y_{(S')}\to X\times_S Y$ es fielmente plano (2.2.13).
Los elementos de $\operatorname{Hom}_S(X,Y)$ corresponden biunívocamente a
los subpreesquemas de $X\times_S Y$ que son los \emph{grafos} de estos
$S$-morfismos (\textbf I, 5.3.11), y, si $Z_f$ es el grafo de
$f\in\operatorname{Hom}_S(X,Y)$, se tiene $Z_{f_{(S')}}=g^{-1}(Z_f)$
(\textbf I, 5.3.12). Basta, pues, aplicar a $g$ la proposición (2.2.15).

\begin{proposition}[2.2.17]
\label{IV.2.2.17-fr}
Sean $A$ un anillo y $B$ una $A$-álgebra tal que $B$ sea un $A$-módulo
fielmente plano y de presentación finita. Entonces el homomorfismo estructural
$\varphi:A\to B$ es un isomorfismo del $A$-módulo $A$ con un sumando
directo del $A$-módulo $B$. Si $A$ es un anillo local, $B$ es un $A$-módulo
libre y existe una base de este módulo que contiene el elemento unidad de $B$.
\end{proposition}

\oldpage[IV-2]{14}
En virtud de Bourbaki, \emph{Alg. comm.}, cap.~II, \S~3,
núm.~3, prop.~12, basta demostrar la proposición cuando $A$
es un anillo local; se sabe entonces (\emph{loc. cit.},
núm.~2, cor.~2 de la prop.~5) que $B$ es un $A$-módulo libre
de tipo finito, y la conclusión resulta de \emph{loc. cit.}, prop.~5.
